\chapter{Extension}
\label{ch4}

Section~\ref{sec:indistinct} attributed the indistinctness of the samples in Chapter~\ref{ch3} to
two causes: a training budget roughly three orders of magnitude below the published configuration,
and the tendency of a squared-error objective to regress towards the mean wherever the network is
uncertain. The Conclusions proposed remedies for both. This chapter carries them out.

The question changes accordingly. Chapter~\ref{ch3} asked which of three denoising networks
generates the best images when all are trained identically; this chapter holds the proposed ResUNet
and the diffusion framework fixed and asks how much of the gap to the published result can be closed
by changing only how the network is trained and sampled. The changes are applied one at a time as a
cumulative ladder, so that the contribution of each can be read from the difference between adjacent
rungs.

% 결과는 주말 72시간 실행(DissertationExtension.ipynb, PRESET = "full")이 끝난 뒤
% model/extension/extension_study.json 에서 4.7절 표로 옮긴다
Sections~\ref{sec:ext-motivation} to~\ref{sec:ext-budget} set out the design, the reasoning behind
every change and the computational budget. Section~\ref{sec:ext-results} states in advance how the
results are to be read and reports them once the runs are complete.

\section{Motivation}
\label{sec:ext-motivation}

The two causes identified in Chapter~\ref{ch3} call for different remedies, and it is worth stating
each precisely.

The first is scale. The networks of Chapter~\ref{ch3} have 2.47 million parameters and were trained
for 11\,700 optimisation steps, against 35.7 million parameters and 800\,000 steps for the CIFAR-10
configuration of Ho et al.\ \parencite{arXiv:2006.11239} (Table~\ref{tab:budget}). Of the two
ratios, the one in steps is the larger by a factor of five, and so is the first target here.

The second is a property of the objective. The network is trained to minimise
\begin{equation}
    \mathcal{L}_{\text{simple}} = \mathbb{E}_{t,\,x_0,\,\varepsilon}
    \left[\left\lVert \varepsilon - \varepsilon_\theta(x_t, t) \right\rVert^2\right],
    \qquad x_t = \sqrt{\bar\alpha_t}\,x_0 + \sqrt{1-\bar\alpha_t}\,\varepsilon,
    \label{eq:ext-simple-loss}
\end{equation}
and the function that minimises a squared error is the conditional expectation of its target,
\begin{equation}
    \varepsilon_\theta^{*}(x_t, t) = \mathbb{E}\left[\varepsilon \mid x_t\right]
    \quad\Longrightarrow\quad
    \hat{x}_0 = \frac{x_t - \sqrt{1-\bar\alpha_t}\,\varepsilon_\theta^{*}(x_t, t)}{\sqrt{\bar\alpha_t}}
    = \mathbb{E}\left[x_0 \mid x_t\right].
    \label{eq:ext-cond-mean}
\end{equation}
Where several clean images are consistent with a noisy one, the ideal network therefore predicts
their average, and the average of several plausible images is smoother than any of them. The
reverse process counteracts this by committing to one outcome a little at a time, as argued in
Section~\ref{sec:indistinct}, but only if each of its thousand predictions is accurate. Anything
that reduces the error of an individual prediction, including noise in the weights themselves,
therefore also addresses the second cause.

\section{Experimental Design}
\label{sec:ext-design}

\subsection{What Is Held Fixed}

Everything that defines the generative model is unchanged from Chapter~\ref{ch3}: $T = 1000$
diffusion steps under the cosine schedule \parencite{arXiv:2102.09672}, the $\varepsilon$-prediction
objective, the ancestral sampler with the estimate of $x_0$ clipped at every step
(Section~\ref{sec:reverse-note}), the ResUNet family of Section~\ref{sec:simulation}, a batch size
of 128, the Adam optimiser \parencite{arXiv:1412.6980} at a learning rate of $2 \times 10^{-4}$,
random horizontal flipping, seed~0, and a Fr\'echet Inception Distance computed from 2048 samples.
Every sample in every rung is drawn from the same initial noise.

\subsection{The Ladder}

Table~\ref{tab:ladder} lists the seven rungs. Each inherits every setting of the one above it and
changes one thing, so the difference in FID between adjacent rungs,
\begin{equation}
    \Delta\mathrm{FID}_i = \mathrm{FID}(E_i) - \mathrm{FID}(E_{i-1}),
    \label{eq:ext-delta}
\end{equation}
is attributable to that change. Two qualifications apply. Three rungs (E2, E5 and E6) change more
than one setting at once, for reasons given below, and their contributions cannot be divided further.
And because the ladder is cumulative, each contribution is measured in the presence of all the
changes that precede it; applied in a different order, a change could contribute a different amount.

\begin{table}[htbp]
    \centering
    \caption{The cumulative ladder. Each rung adds the change in its row to all of those above it.
    Widths are listed for CIFAR-10; the corresponding CelebA widths are $(32, 64, 128)$ and
    $(96, 192, 256)$.}
    \label{tab:ladder}
    \small
    \begin{tabular}{llrlc}
        \toprule
        Rung & Change added & Epochs & Widths & Blocks per stage \\
        \midrule
        E0 & none (Chapter~\ref{ch3}, seed 0)            & 30  & $(64, 128)$ & 1 \\
        E1 & moving average of the weights              & 30  & $(64, 128)$ & 1 \\
        E2 & warm-up, gradient clipping, dropout        & 30  & $(64, 128)$ & 1 \\
        E3 & min-SNR loss weighting, $\gamma = 5$         & 30  & $(64, 128)$ & 1 \\
        E4 & five times the training                      & 150 & $(64, 128)$ & 1 \\
        E5 & increased capacity                           & 150 & $(128, 256, 256)$ & 2 \\
        E6 & long training, EMA decay 0.9999              & 2000\textsuperscript{a} &
        $(128, 256, 256)$ & 2 \\
        \bottomrule
        \multicolumn{5}{l}{\footnotesize \textsuperscript{a} 700 on CelebA; see
        Section~\ref{sec:ext-long}.}
    \end{tabular}
\end{table}

The complete ladder is run on CIFAR-10, which is where the contribution of each change is isolated.
On CelebA only E0, E5 and E6 are run. The purpose there is to establish whether the overall effect
carries over to a second dataset at twice the resolution, and the difference between the first rung
and the last is sufficient for that.

\subsection{Validation of the Starting Point}

E0 repeats the seed-0 ResUNet run of Chapter~\ref{ch3} exactly, and serves as a check on the new
pipeline: if it did not reproduce the earlier result, no other row could be trusted. The criterion
adopted is agreement within one seed-to-seed standard deviation of that result, which
Table~\ref{tab:fid} puts at $21.13$ for this network on CIFAR-10. A preliminary evaluation of the E0
checkpoint gives an FID of $89.57$ against the $94.93$ of Chapter~\ref{ch3}, a difference of $5.36$,
well inside that criterion.

\section{Rationale for Each Change}
\label{sec:ext-rationale}

\subsection{Exponential Moving Average of the Weights (E1)}
\label{sec:ext-ema}

Stochastic gradient descent at a constant learning rate does not settle at a minimum; it fluctuates
around one, with an amplitude governed by the learning rate and the variance of the gradient. The
last iterate is one draw from that fluctuation. Polyak and Juditsky showed that the average of the
iterates is a better estimate of the minimum than any single one of them, and converges at the
optimal rate \parencite{polyak1992averaging}. Diffusion models apply the same idea as an exponential
moving average,
\begin{equation}
    \bar\theta_k = d_k\,\bar\theta_{k-1} + (1 - d_k)\,\theta_k,
    \qquad d_k = \min\!\left(d,\; \frac{1 + k}{10 + k}\right),
    \label{eq:ext-ema}
\end{equation}
where $\theta_k$ are the weights after step $k$ and $\bar\theta_k$ the averaged copy. Training
continues on $\theta_k$; only sampling uses $\bar\theta_k$. The practice is standard in the
diffusion literature: Ho et al.\ sample from an average with $d = 0.9999$
\parencite{arXiv:2006.11239}, and Song and Ermon identify it as one of the techniques on which the
stability of score-based models depends \parencite{arXiv:2006.09011}.

For a constant decay the average places weight $(1-d)\,d^{\,k-j}$ on the iterate $j$ steps in the
past, so it effectively averages over the last $1/(1-d)$ steps. The warm-up in
Equation~\eqref{eq:ext-ema} exists because a constant decay misbehaves in a short run. With
$d = 0.9999$ and no warm-up, a quantity $d^{\,k}$ of the average still belongs to the random
initialisation after $k$ steps, which is $31\%$ at the end of a 30-epoch run of 11\,700 steps. Under
the warm-up the effective window grows as $(10 + k)/9$, so the average always covers roughly the most
recent ninth of training, until the window reaches the cap $1/(1-d)$ at step $(10d - 1)/(1 - d)$.
Figure~\ref{fig:ext-ema-window} shows the resulting window against the length of each rung.

\begin{figure}[htbp]
    \centering
    \includegraphics[width=0.8\linewidth]{ext_ema_window.png}
    \caption{Effective averaging window $1/(1 - d_k)$ of the moving average in
    Equation~\eqref{eq:ext-ema}. The two caps coincide throughout the warm-up, which ends at step
    17\,990 for $d = 0.9995$ and at step 89\,990 for $d = 0.9999$. Dotted lines mark the length of
    each rung. Drawn from Equation~\eqref{eq:ext-ema}.}
    \label{fig:ext-ema-window}
\end{figure}

Two consequences follow. The first is that the thirty-epoch rungs E1 to E3 end at step 11\,700,
before either cap is reached, so their average is determined by the warm-up alone and the choice of
cap has no effect on them. The second is that the cap matters only for the long rungs, and there it
sets a trade-off between variance and lag. A longer window averages out more of the fluctuation but
also trails further behind weights that are still improving. E4 and E5 use $d = 0.9995$, whose window
of 2000 steps is about five epochs and $3.4\%$ of the run. E6 uses $d = 0.9999$, the value of Ho et
al., whose window of 10\,000 steps is $1.3\%$ of its 780\,000; late in a run of that length the
weights drift slowly enough that the longer average costs little in lag. The change of cap is part
of E6 rather than a rung of its own for this reason, since it is a consequence of the length of
training rather than an independent choice.

\subsection{Learning-Rate Warm-Up, Gradient Clipping and Dropout (E2)}
\label{sec:ext-stabilisers}

The three changes in E2 share a purpose. Each makes training more robust without altering what is
being learned, each is used in the reference implementation of Ho et al., and each is expected to
have a small effect individually at thirty epochs. Giving each a rung of its own would have added two
thirty-epoch runs to measure differences that a single seed could not resolve.

\paragraph{Warm-up.} Adam scales each step by a running estimate of the second moment of the
gradient,
\begin{equation}
    \theta_{k+1} = \theta_k - \eta_k\,\frac{\hat m_k}{\sqrt{\hat v_k} + \epsilon},
    \label{eq:ext-adam}
\end{equation}
where $\hat m_k$ and $\hat v_k$ are the bias-corrected first and second moment estimates
\parencite{arXiv:1412.6980}. In the first steps $\hat v_k$ is estimated from very few gradients, and
Liu et al.\ show that the variance of the resulting adaptive learning rate is then large enough to
drive training towards poor regions of the loss surface \parencite{arXiv:1908.03265}. A linear warm-up,
\begin{equation}
    \eta_k = \eta\,\min\!\left(1,\; \frac{k+1}{K}\right),
    \label{eq:ext-warmup}
\end{equation}
keeps the step small until the estimate has settled. The same device is used to stabilise training
at large batch sizes \parencite{arXiv:1706.02677}. The reference implementation uses $K = 5000$; a
value of $K = 2000$ is used here because the thirty-epoch rungs last only 11\,700 steps, and 5000
would spend $43\%$ of them warming up against $17\%$ for 2000.

\paragraph{Gradient clipping.} The training loss mixes timesteps of very different difficulty, and
an occasional batch produces a gradient far larger than typical. Rescaling the gradient whenever its
norm exceeds a threshold $c$,
\begin{equation}
    g \leftarrow g \cdot \min\!\left(1,\; \frac{c}{\lVert g \rVert_2}\right),
    \label{eq:ext-clip}
\end{equation}
bounds the size of the step while preserving its direction. The technique was introduced to control
exploding gradients in recurrent networks \parencite{arXiv:1211.5063}; $c = 1$ is the value of the
reference implementation, and the norm is taken over all parameters together.

\paragraph{Dropout.} Rungs E4 to E6 pass over the same fifty thousand images between 150 and 2000
times, and a network trained that long can begin to fit its training set rather than the
distribution behind it. Dropout counters this by zeroing each activation independently with
probability $p$ during training and rescaling the survivors,
\begin{equation}
    \tilde h = \frac{m \odot h}{1 - p}, \qquad m_i \sim \mathrm{Bernoulli}(1 - p),
    \label{eq:ext-dropout}
\end{equation}
so that no unit can rely on the presence of any other \parencite{srivastava2014dropout}. Ho et al.\
use $p = 0.1$ on CIFAR-10, placed before the second convolution of each residual block, and the same
value and placement are used here. Dropout must be disabled when sampling. Left active, it would
inject independent noise into each of the thousand predictions of the reverse process, and every
sampling and evaluation routine therefore switches the network to inference mode explicitly.

\subsection{Min-SNR Loss Weighting (E3)}
\label{sec:ext-minsnr}

The signal-to-noise ratio of the noisy image at timestep $t$ is \parencite{arXiv:2107.00630}
\begin{equation}
    \mathrm{SNR}(t) = \frac{\bar\alpha_t}{1 - \bar\alpha_t}.
    \label{eq:ext-snr}
\end{equation}
Substituting the estimate of $x_0$ from Equation~\eqref{eq:ext-cond-mean} shows that the error in
the predicted noise and the error in the implied clean image are related by
\begin{equation}
    x_0 - \hat{x}_0 = \sqrt{\frac{1 - \bar\alpha_t}{\bar\alpha_t}}\,
    \left(\varepsilon_\theta - \varepsilon\right)
    \quad\Longrightarrow\quad
    \left\lVert \varepsilon - \varepsilon_\theta \right\rVert^2
    = \mathrm{SNR}(t)\,\left\lVert x_0 - \hat{x}_0 \right\rVert^2 .
    \label{eq:ext-eps-x0}
\end{equation}
The standard objective is therefore an $x_0$ objective in which timestep $t$ is weighted by
$\mathrm{SNR}(t)$. Under the cosine schedule that weight is $24\,221$ at $t = 0$ and $2.4 \times
10^{-9}$ at $t = 999$ (Figure~\ref{fig:ext-minsnr}a), so the objective devotes most of its accuracy
to the nearly clean timesteps, where denoising is easiest and matters least for the sample. Hang et
al.\ describe the timesteps as competing tasks whose gradients conflict, and propose capping the
$x_0$ weight at a constant $\gamma$ \parencite{arXiv:2303.09556}. In the $\varepsilon$
parameterisation used here this becomes
\begin{equation}
    \mathcal{L}_{\text{min-SNR}} = \mathbb{E}_{t,\,x_0,\,\varepsilon}
    \left[ w(t)\,\left\lVert \varepsilon - \varepsilon_\theta(x_t, t) \right\rVert^2 \right],
    \qquad
    w(t) = \frac{\min\{\mathrm{SNR}(t),\, \gamma\}}{\mathrm{SNR}(t)} .
    \label{eq:ext-minsnr}
\end{equation}
As $\gamma \to \infty$ the weight tends to one and the standard objective is recovered. With the
authors' value of $\gamma = 5$, the weight falls below one wherever $\mathrm{SNR}(t) > 5$, which
under this schedule is the first $26.1\%$ of timesteps, $t \le 260$
(Figure~\ref{fig:ext-minsnr}b). Its mean over all timesteps is $0.824$, and $18.8\%$ of timesteps
receive less than half their former weight. Because Adam normalises each step by the magnitude of
the gradient, the reduction in the overall scale of the loss does not by itself change the step
size; what changes is the relative emphasis among timesteps.

\begin{figure}[htbp]
    \centering
    \includegraphics[width=\linewidth]{ext_minsnr_weight.png}
    \caption{(a) Signal-to-noise ratio of the noisy image under the clipped cosine schedule, with the
    cap $\gamma = 5$ dashed. (b) The resulting weight on the $\varepsilon$-prediction loss. Timesteps
    to the right of $t = 260$ are unaffected. Drawn from Equations~\eqref{eq:ext-snr}
    and~\eqref{eq:ext-minsnr}, following Hang et al.\ \parencite{arXiv:2303.09556}.}
    \label{fig:ext-minsnr}
\end{figure}

Hang et al.\ report that the weighting reaches a given FID several times sooner than the standard
objective, but their evidence comes from networks considerably larger than the 2.47
million-parameter network of rung E3. Its
effect at this scale is therefore not assumed; the rung measures it, and if the weighting does not
help, the result will show it.

\subsection{Five Times the Training (E4)}

E4 extends training from 30 to 150 epochs, or from 11\,700 to 58\,500 steps, and so reduces the
shortfall in optimisation steps against Ho et al.\ from a factor of 68 to one of 14. It is the most
direct response to the first cause and involves no change of method at all.

\subsection{Increased Capacity (E5)}

E5 widens the network and deepens each stage. On CIFAR-10 the widths change from $(64, 128)$ to
$(128, 256, 256)$ with two residual blocks per stage instead of one, which takes the parameter count
from 2.47 million to 31.46 million, $88\%$ of the 35.7 million of Ho et al. The widths and block
count follow the published configuration, whose base width is 128 with two residual blocks at each
resolution. On CelebA, which already has three stages, the widths change from $(32, 64, 128)$ to
$(96, 192, 256)$, taking the count from 2.61 million to 24.41 million.

E5 is trained for the same 150 epochs as E4, so the difference between them is due to capacity
alone. That capacity arrives in three forms at once on CIFAR-10, however: wider layers, a second
block per stage, and a third resolution stage. The rung therefore cannot serve as a test of the depth
hypothesis raised in Section~\ref{sec:distribution}, which requires the number of stages to change
while the width is held fixed.

\subsection{Long Training (E6)}
\label{sec:ext-long}

E6 spends the remainder of the available time on the E5 network. It resumes from the E5 checkpoint,
weights, optimiser state and moving average included, and continues to 2000 epochs on CIFAR-10 and
700 on CelebA. On CIFAR-10 this is 780\,000 optimisation steps, $97.5\%$ of the 800\,000 of Ho et
al., so the shortfall in steps that dominated the budget of Chapter~\ref{ch3} is essentially removed
(Table~\ref{tab:ext-budget}). A CelebA step costs 2.8 times as much as a CIFAR-10 step, and 700 epochs,
or 273\,000 steps, gives the two datasets roughly equal shares of the time.

\begin{table}[htbp]
    \centering
    \caption{Training budget of the extended rungs on CIFAR-10, continuing
    Table~\ref{tab:budget}. E6 removes the shortfall in optimisation steps; the remaining
    difference from the published configuration lies in the architecture.}
    \label{tab:ext-budget}
    \begin{tabular}{lrrrrr}
        \toprule
        & Chapter~\ref{ch3} & E4 & E5 & E6 & Ho et al. \\
        \midrule
        Parameters         & 2.47\,M & 2.47\,M & 31.46\,M & 31.46\,M & 35.7\,M \\
        Optimisation steps & 11\,700 & 58\,500 & 58\,500  & 780\,000 & 800\,000 \\
        Steps relative to Ho et al. & $1/68$ & $1/14$ & $1/14$ & $0.98$ & $1$ \\
        \bottomrule
    \end{tabular}
\end{table}

Resuming rather than starting afresh means that E5 is the 150-epoch point on the trajectory of E6,
not an independent run, and costs nothing beyond its evaluation. It also means that the difference
between them combines the length of training with the change of averaging cap discussed in
Section~\ref{sec:ext-ema}. What remains between E6 and the published configuration is the
architecture: the reference network has four resolution stages rather than three, self-attention at
$16 \times 16$, and a skip connection from every block (Section~\ref{sec:simulation}). Those
differences are deliberately left in place, since this chapter asks what training alone can achieve.

\section{Sampling}
\label{sec:ext-sampling}

Every rung is evaluated with the thousand-step ancestral sampler of Chapter~\ref{ch3}, so that E0
can be compared with the earlier result and every rung with every other. The final rung on each
dataset is evaluated a second time with the deterministic sampler of Song et al.\
\parencite{arXiv:2010.02502}. It visits a subsequence $\tau_1 > \tau_2 > \dots > \tau_S$ of the
timesteps and, with its stochasticity parameter set to zero, updates
\begin{equation}
    x_{\tau_{i+1}} = \sqrt{\bar\alpha_{\tau_{i+1}}}\;\hat{x}_0
    + \sqrt{1 - \bar\alpha_{\tau_{i+1}}}\;\varepsilon_\theta(x_{\tau_i}, \tau_i),
    \qquad
    \hat{x}_0 = \frac{x_{\tau_i} - \sqrt{1 - \bar\alpha_{\tau_i}}\,
    \varepsilon_\theta(x_{\tau_i}, \tau_i)}{\sqrt{\bar\alpha_{\tau_i}}},
    \label{eq:ext-ddim}
\end{equation}
with $\hat{x}_0$ clipped to $[-1, 1]$ as in the ancestral sampler and $\bar\alpha$ taken as one after
the last step, so that the final update returns $\hat{x}_0$ itself. With $S = 100$ evenly spaced
timesteps it is ten times cheaper than the ancestral sampler. It is reported as a separate row because
a change of sampler changes the FID, and mixing the two would break the comparison along the ladder.

\section{Computational Budget}
\label{sec:ext-budget}

The rungs were fitted to a fixed allowance of 72 hours on the machine of Table~\ref{tab:environment}.
Two engineering measures made that allowance go considerably further than the settings of
Chapter~\ref{ch3} would have allowed, and one of them proved to be a prerequisite rather than an
optimisation.

\subsection{Graph Compilation}

The network was compiled with the graph compiler of PyTorch~2 \parencite{ansel2024pytorch}, which
fuses sequences of elementwise operations such as group normalisation and activation into single
kernels. Like the mixed precision and memory layout of Chapter~\ref{ch3}, this changes the arithmetic
used to evaluate the model and not the model itself. Table~\ref{tab:ext-steptime} reports the
measured effect. On the three smaller networks compilation shortens each step by a factor of 1.6 to
2.0. A second measure, removing the transfer of the loss from the graphics card to the host at every
step so that the processor no longer waits for each step to finish before issuing the next, gained a
further $8\%$ on the smallest network and raised the utilisation of the card from $91\%$ to $97\%$.

\begin{table}[htbp]
    \centering
    \caption{Measured time per optimisation step at a batch size of 128, with the card otherwise
    idle. Eager execution corresponds to the settings of Chapter~\ref{ch3}; its value for the
    ResUNet on CIFAR-10 reproduces the $16.0$\,s per epoch of Table~\ref{tab:cost}.}
    \label{tab:ext-steptime}
    \begin{tabular}{llrrr}
        \toprule
        Dataset & Network & Eager & Compiled & Speed-up \\
        \midrule
        CIFAR-10 & 2.47\,M (E0--E4)  & $41.0$\,ms  & $21.7$\,ms  & $1.9\times$ \\
        CIFAR-10 & 31.46\,M (E5, E6) & $221$\,ms   & $141$\,ms   & $1.6\times$ \\
        CelebA   & 2.61\,M (E0)      & $84.1$\,ms  & $43.1$\,ms  & $2.0\times$ \\
        CelebA   & 24.41\,M (E5, E6) & $3465$\,ms  & $390$\,ms   & $8.9\times$ \\
        \bottomrule
    \end{tabular}
\end{table}

The last row is not a speed-up in the ordinary sense. Executed eagerly, the large CelebA network
requires $12.9$\,GB at its peak, more than the 12\,GB the card provides. The graphics driver does not
report this as an error but silently borrows system memory, and each step then spends most of its
time waiting on transfers: the step takes $3.5$ seconds and the power drawn by the card falls to 88\,W,
against 190\,W when it is computing. Compilation reduces the peak to $7.2$\,GB by not materialising
intermediate tensors that the fused kernels no longer need, and the network fits. Without it, the
700 epochs of E5 and E6 on CelebA would have taken more than 260 hours. The implementation caps the memory available to
the process at $90\%$ of the card, so that an overflow of this kind stops the run with an error
rather than slowing it unnoticed.

\subsection{Compiled Sampling}

The thousand steps of the reverse process were compiled in the same way and, in addition, recorded
as a replayable sequence of kernels, which removes the cost of launching each kernel individually.
This shortens sampling by a factor of 2.0 to 3.0, and the 2048 samples needed for the FID of the
small CIFAR-10 network take $1.5$ minutes instead of $3.8$. Compilation reorders floating-point
operations, so the samples are not bit-for-bit identical to those of eager execution even from the
same noise. Its effect on the measurement was checked directly on the E0 checkpoint: eager sampling
gives an FID of $89.57$ and an Inception Score of $4.649$, compiled sampling $90.97$ and $4.650$. The
difference of $1.4$ FID is less than a fifth of the smallest seed-to-seed standard deviation observed
in Chapter~\ref{ch3}, and since every rung is evaluated the same way it cannot affect
$\Delta\mathrm{FID}$.

\subsection{Schedule}

Table~\ref{tab:ext-schedule} gives the resulting schedule. Training accounts for 61 of the 63 hours,
and the four large-network rungs for $96\%$ of the total. The remaining nine hours of the allowance
are held in reserve against the slowing of a card run continuously for three days and against
interruptions; checkpoints are written after every epoch, so an interrupted run resumes where it
stopped.

\begin{table}[htbp]
    \centering
    \caption{Schedule of the complete ladder, computed from the measured step times of
    Table~\ref{tab:ext-steptime}. One epoch is 390 steps on both datasets.}
    \label{tab:ext-schedule}
    \begin{tabular}{lrrr}
        \toprule
        Item & Epochs & Seconds per epoch & Time \\
        \midrule
        CIFAR-10 E0--E3              & $4 \times 30$ & 8.5  & 17\,min \\
        CIFAR-10 E4                  & 150           & 8.5  & 21\,min \\
        CIFAR-10 E5                  & 150           & 55   & 2\,h 18\,min \\
        CIFAR-10 E6 (from E5)        & $+1850$       & 55   & 28\,h 16\,min \\
        CelebA E0                    & 30            & 16.8 & 8\,min \\
        CelebA E5                    & 150           & 152  & 6\,h 20\,min \\
        CelebA E6 (from E5)          & $+550$        & 152  & 23\,h 13\,min \\
        Compilation, once per rung   &               &      & 8\,min \\
        \midrule
        Training                     &               &      & 61\,h \\
        Evaluation and figures       &               &      & 1\,h 40\,min \\
        \midrule
        Total                        &               &      & 63\,h \\
        \bottomrule
    \end{tabular}
\end{table}

\section{Results}
\label{sec:ext-results}

\subsection{How the Results Are to Be Read}

The rules for interpreting the ladder are fixed here, before the results are known, so that they
cannot be chosen to suit them.

\begin{enumerate}
    \item \textbf{Only $\Delta\mathrm{FID}$ supports a causal claim.} The absolute FID of a rung
    includes the effect of every change above it.

    \item \textbf{Each rung is a single run.} Chapter~\ref{ch3} measured the seed-to-seed standard
    deviation of FID for this network at $21.13$ on CIFAR-10 and $23.98$ on CelebA. A rung whose
    $\Delta\mathrm{FID}$ is smaller in magnitude than that cannot be said to have had an effect, and
    will be reported as inconclusive rather than as a small improvement.

    \item \textbf{E0 must reproduce Chapter~\ref{ch3}.} The criterion is agreement with $94.93$ to
    within one standard deviation. The preliminary evaluation reported in Section~\ref{sec:ext-design}
    already meets it.

    \item \textbf{Bundled rungs are reported as bundles.} E2 is one result for three stabilisers,
    E5 one result for width, depth and a third stage on CIFAR-10, and E6 one result for training
    length and averaging cap together.
\end{enumerate}

\subsection{Measurements}

% 실행이 끝나면 model/extension/extension_study.json 에서 옮긴다.
% 키: "<데이터셋>|<단계>|ancestral" 과 "<데이터셋>|E6|ddim100"
% loss_final -> 손실, fid -> FID, is_mean -> IS. ΔFID 는 바로 위 행과의 차
\noindent\textit{The measurements below are completed from the output of the 72-hour run described
in Section~\ref{sec:ext-budget}. Epochs and parameter counts are fixed by the design and are given
in advance.}

\begin{table}[htbp]
    \centering
    \caption{Results of the ladder. $\Delta$FID is the difference from the row above and is the
    contribution of that rung. The final row of each dataset uses the deterministic sampler and is
    compared with the row above it.}
    \label{tab:ext-results}
    \small
    \begin{tabular}{llrrrrrr}
        \toprule
        Dataset & Rung & Epochs & Parameters & Final loss & FID & $\Delta$FID & IS \\
        \midrule
        CIFAR-10 & E0 & 30   & 2\,473\,475  & --- & --- & ---  & --- \\
        CIFAR-10 & E1 & 30   & 2\,473\,475  & --- & --- & ---  & --- \\
        CIFAR-10 & E2 & 30   & 2\,473\,475  & --- & --- & ---  & --- \\
        CIFAR-10 & E3 & 30   & 2\,473\,475  & --- & --- & ---  & --- \\
        CIFAR-10 & E4 & 150  & 2\,473\,475  & --- & --- & ---  & --- \\
        CIFAR-10 & E5 & 150  & 31\,460\,739 & --- & --- & ---  & --- \\
        CIFAR-10 & E6 & 2000 & 31\,460\,739 & --- & --- & ---  & --- \\
        CIFAR-10 & E6, deterministic, 100 steps & 2000 & 31\,460\,739 & & --- & --- & --- \\
        \midrule
        CelebA   & E0 & 30   & 2\,612\,419  & --- & --- & ---  & --- \\
        CelebA   & E5 & 150  & 24\,408\,291 & --- & --- & ---  & --- \\
        CelebA   & E6 & 700  & 24\,408\,291 & --- & --- & ---  & --- \\
        CelebA   & E6, deterministic, 100 steps & 700 & 24\,408\,291 & & --- & --- & --- \\
        \bottomrule
    \end{tabular}
\end{table}

\begin{figure}[htbp]
    \centering
    \IfFileExists{imgs/ext_fid_ladder.png}{%
        \includegraphics[width=\linewidth]{ext_fid_ladder.png}%
    }{%
        \fbox{\parbox[c][4cm][c]{0.9\linewidth}{\centering\itshape
        FID by rung; generated by Section 10 of the extension notebook}}%
    }
    \caption{FID at each rung, with the seed-0 ResUNet result of Chapter~\ref{ch3} dashed.}
    \label{fig:ext-fid-ladder}
\end{figure}

\begin{figure}[htbp]
    \centering
    \IfFileExists{imgs/ext_samples_cifar10.png}{%
        \includegraphics[width=\linewidth]{ext_samples_cifar10.png}%
    }{%
        \fbox{\parbox[c][3cm][c]{0.9\linewidth}{\centering\itshape
        CIFAR-10 samples by rung; generated by Section 10 of the extension notebook}}%
    }\\[3mm]
    \IfFileExists{imgs/ext_samples_celeba.png}{%
        \includegraphics[width=\linewidth]{ext_samples_celeba.png}%
    }{%
        \fbox{\parbox[c][3cm][c]{0.9\linewidth}{\centering\itshape
        CelebA samples by rung; generated by Section 10 of the extension notebook}}%
    }
    \caption{Samples from each rung on CIFAR-10 (top) and CelebA (bottom). Every panel starts from
    the same initial noise, so differences between panels are due to the training configuration
    alone.}
    \label{fig:ext-samples}
\end{figure}

\section{Limitations}
\label{sec:ext-limitations}

The most serious limitation is the one Chapter~\ref{ch3} was at pains to establish: a single run
cannot distinguish an effect from the seed-to-seed variation unless the effect is large. The
72-hour allowance was spent on length of training rather than on repetition, because the question
of this chapter concerns absolute quality, and the reading rules above are the price of that choice.

The ladder measures each change in the context of those above it. The moving average in particular
is likely to be worth more in a short run, where the final iterate is far from settled, than it would
be if added at the end of a long one; E1 measures it at thirty epochs and makes no claim beyond that.
Three rungs combine several changes, as set out in Section~\ref{sec:ext-rationale}.

The FID is again computed from 2048 samples, for comparability with Chapter~\ref{ch3}, and remains
biased upward at that size; the absolute values should not be compared with published figures computed
from fifty thousand. Compiled sampling introduces a numerical difference that was measured at
$1.4$~FID on one checkpoint.

Several remedies were deliberately left out because they would change the method rather than its
training. Predicting the velocity $v$ instead of the noise alters the implicit weighting across
timesteps and hence the regression to the mean discussed in Section~\ref{sec:ext-motivation}
\parencite{arXiv:2202.00512}. Classifier-free guidance is the most effective single remedy for
indistinct samples in current practice \parencite{arXiv:2207.12598}, but it requires a conditional
model and does not apply to the unconditional generation studied here. And the architectural
differences from the reference network, among them self-attention and a fourth resolution stage,
are left in place so that the result measures what training alone can achieve.

\section{Summary}

This chapter applies to the proposed network the remedies that Chapter~\ref{ch3} and the Conclusions
identified for its indistinct samples. Seven cumulative rungs add an average of the weights, three
training stabilisers, a reweighting of the loss across timesteps, five times the training, nearly
thirteen times the parameters and, finally, sixty-seven times the optimisation steps of the original
experiment, which brings the step count on CIFAR-10 to within $2.5\%$ of the published configuration.
The reasoning behind each change and the value chosen for it are set out in
Section~\ref{sec:ext-rationale}, and the rules by which the results are to be read are fixed in
Section~\ref{sec:ext-results} ahead of them.

Two findings are already in hand. The pipeline reproduces the seed-0 result of Chapter~\ref{ch3}
to within a quarter of a standard deviation, which validates every rung built on it. And graph
compilation is not merely an acceleration at this scale but a condition of the experiment: without
it the largest network does not fit in the memory of the card, and the driver's silent fallback to
system memory makes it nine times slower without reporting any fault.
